\section{Discussion}
The direction and magnitude of the correlations suggest that some measured weather variables co-vary with RSSI, SNR, and throughput in this dataset. However, the predictive experiments explain at most about 21\% of holdout variance, and flexible random-forest/KNN models generally do not improve performance. This pattern is consistent with weak signal, confounding by collection session or time, limited covariates, or train/test dependence; the repository cannot distinguish among these explanations.

The positive association of rain accumulators with link metrics should not be interpreted as rain improving propagation. Accumulator variables may encode time or collection-session identity, and pressure and rain measures may co-vary with unrecorded changes in geometry or hardware. Furthermore, random row splitting allows temporally adjacent, similarly averaged observations to occur in both partitions. The reported scores therefore characterize the stored pipeline, not expected performance on a new storm, site, or radio link.

\section{Limitations}
The principal limitations are: undocumented hardware and measurement units; an unknown data-collection protocol; only one identified location name (Wilson Hall) without geometry; 2,432 duplicate rows in the consolidated weather file; possible overlapping merge windows; no temporal or session-aware validation; no uncertainty intervals or multiple-testing correction; unit-mixing in averaged-feature correlations; redundant rain variables; and failed scalar interpretation of vector-valued particle fields. The scripts also use different feature sets and test fractions across targets, limiting direct model comparison. Generated images do not embed metric values or provenance, and saved console outputs are absent.

\section{Future Work}
Future work should document instruments and units, preserve explicit collection-session identifiers, decide whether repeated weather rows are valid, parse vector-valued particle distributions into justified summaries, and audit timestamp/time-zone alignment. Models should be evaluated with blocked or leave-one-session-out validation and confidence intervals. A physically motivated feature set should replace arithmetic averages across unlike units, and collinearity between cumulative rain fields should be addressed. Additional measurements across locations, distances, radio settings, and weather regimes are required before generalization.
